\documentclass[12pt]{article}
\usepackage[utf8]{inputenc}
\usepackage[T1]{fontenc}
\usepackage{hyperref}
\title{Trinitario Language Resources for K-12 Curriculum Development}
\author{Compiled by Claude Sonnet 5 (claude-sonnet-5) for Andr\'as Kornai}
\date{}
\begin{document}
\maketitle

\section{Introduction}
Trinitario (Mojeño Trinitario; ISO 639-3 trn) is an endangered Arawakan language of the Bolivian lowlands, spoken around Trinidad and San Lorenzo and San Francisco de Moxos. Searches turned up scholarly description and archive material but no school materials. Sources disagree on whether Trinitario and Ignaciano are separate languages or dialects of one language.

\section{Primary Educational Resources}
No primers, textbooks or school programmes for Trinitario were found in the search results.

\section{Linguistic and Reference Materials}
\begin{itemize}
\item Fran\c{c}oise Rose, ``Mojeño Trinitario'' (Illustrations of the IPA), Journal of the International Phonetic Association: \url{https://www.cambridge.org/core/journals/journal-of-the-international-phonetic-association/article/abs/mojeno-trinitario/8BD8F413069BDC73EE0906A27BF755D8}
\item Profile of Trinitario at the Endangered Languages Project: \url{https://www.endangeredlanguages.com/lang/2910?hl=en}
\item ASJP entry for Trinitario (word list): \url{https://asjp.clld.org/languages/TRINITARIO}
\item The searches also mention \emph{Gram\'atica Moje\~na Trinitaria} (Ib\'a\~nez Noza et al., 2007 and 2009) and a Trinitario dictionary in the Intercontinental Dictionary Series. I did not see a URL for either, so neither is cited with a link.
\end{itemize}

\section{Religious and Community Resources}
None were found in the search results.

\section{Digital and Audio Resources}
\begin{itemize}
\item Profile of Fran\c{c}oise Rose, who has described Trinitario and plans a grammar: \url{https://blogs.uoregon.edu/uolinguistics/?p=133}
\end{itemize}
The searches mention a DoReCo corpus of Mojeño Trinitario audio compiled by Rose. I did not see its URL.

\section{Curriculum Development Context}
Little exists for Trinitario. The reference description says it is not known to be taught in schools, and no curriculum materials were found. A curriculum would have to be built from scratch, drawing on the linguistic descriptions above and working with speaker communities.
\end{document}